%% ----------------------------------------------------------------
%% Optimisation problem sheet (lectures: optimisation, sgd, constrainedOpt)
%% ----------------------------------------------------------------
\documentclass{sotonExamBoxes}    % use option [answers] to include them
%% ----------------------------------------------------------------
\usepackage{graphicx}
\usepackage{alltt}
\usepackage{bm}
\usepackage{amssymb}
\usepackage{amsfonts}
\usepackage{amsmath}
\thinmuskip=1mu
%% macros.tex -- shared macro definitions for the problem sheets.
%%
%% Every sheet inputs this instead of carrying its own copy. Input it after
%% the \usepackage lines: the definitions below use bm, amsmath and amssymb.
%%
%% Six macros were defined inconsistently across the sheets when this file
%% was made; the choices are marked CONFLICT below.




\usepackage{stmaryrd}

\renewcommand{\ttdefault}{pcr}

\newcommand{\tr}{\textsf{T}}

\newcommand{\e}[1]{{\rm e}^{#1}}

\newcommand{\E}[1]{{\rm e}^{{\displaystyle #1}}}

\newcommand{\logg}[1]{\log\left(#1\right)}

\newcommand{\Prob}[1]{\mathbb{P}\left(#1\right)}

\newcommand{\Step}{\mathop{\mathrm{O}\hspace{-15pt}\rule[8.5pt]{11pt}{1.2pt}\hspace{2pt}}}

\newcommand\diag{\mathop{\mathrm{diag}}}

\newcommand{\grad}{\bm{\nabla}}

\newcommand{\dd}{\mathrm{d}}

\newcommand{\class}{\mathcal{C}}

\newcommand{\data}{\mathcal{D}}

%% CONFLICT: gaussianProcess.tex had {U}{eur}{b}{n} (Euler bold);
%% every other sheet had this. Matrices in that one sheet now look
%% different -- see the report.

\DeclareMathAlphabet{\mat}{OT1}{cmss}{bx}{n}

%% The identity needs to be distinguishable. In cmss a bold I is a plain
%% vertical bar with no crosspieces, so \mat{I} reads as a rule rather than a
%% letter. Euler's bold I keeps its serifs, so the identity alone is set in
%% Euler while everything else stays cmss.
\DeclareMathAlphabet{\matEuler}{U}{eur}{b}{n}
\newcommand{\Id}{\matEuler{I}}

%% CONFLICT and BUG FIX: 16 sheets defined this using \qpartial,
%% which is defined nowhere, so \pd was broken in all of them.
%% Only problem2.tex had the working \partial form, kept here.

\newcommand{\pd}[2]{\ensuremath{\frac{\partial #1}{\partial #2}}\xspace}

\newcommand{\M}[1]{\ensuremath{\boldsymbol{#1}}\xspace}

\newcommand{\bst}{\rule{0mm}{5mm}}

%% CONFLICT: some sheets omitted the \! . Kept the tighter spacing,
%% which is what every sheet that actually uses \av had.

\newcommand{\av}[2][\,]{\mathbb{E}_{#1}\!\left[ {\strut #2} \right]}

\newcommand{\len}[1]{\| #1 \|}

\newcommand{\inner}[2]{\left\langle #1, #2\right\rangle}

\newcommand{\squeeze}{\setlength{\itemsep}{-2pt}}

%% CONFLICT: two spacings, identical output in maths mode.

\newcommand{\bra}[1]{\left( #1 \right)}

\newcommand{\pred}[1]{\left\llbracket { \small #1} \right\rrbracket}

\newcommand{\mln}{\hat{f}_n(\bm{x})}


\newcommand{\Tr}{\mathop{\mathrm{tr}}\,}

%% NOT SHARED: gaussianProcess.tex redefines \det to take an argument
%% (\det{\mat{V}}), while calculus.tex and vectorSpaces.tex use the standard
%% no-argument operator (\det(\mat{X})). The two cannot coexist, so that
%% redefinition stays local to gaussianProcess.tex.

%% NOT SHARED: \ml and \mm mean different things in different sheets.
%% amgm and convexity use a no-argument \ml and \mm = \hat{m}_c(x);
%% ensembleLearning uses a one-argument \ml and \mm = \hat{m}(x).
%% Both forms are in use, so each sheet keeps its own.

\graphicspath{{figures/}}
\usepackage{parskip}
\AddToHook{env/samepage/begin}{%
\AddToHook{para/before}{\nopagebreak}%
\AddToHook{para/after}{\nopagebreak}%
}
\begin{document}
\unitTitle{Optimisation Problem Sheet}
\authors{Adam Pr\"ugel-Bennett}
\unitcode{AICE3001/AICE6002}
\semester{2}
\year{2022}
\duration{45}{0}

\maketitle

This sheet covers the three optimisation lectures.  Question~1 looks
at gradient descent and Newton's method on simple functions, where we
can see exactly what happens.  Questions~2 and~3 use Lagrange
multipliers and duality.  Question~4 asks you to check some of these
results in Python.  The early parts of each question are short; the
later parts ask you to put the pieces together.

% ***************************** question 1 **************************

\begin{question}{20}
  Gradient descent with learning rate $r$ updates
  $w \leftarrow w - r\,f'(w)$ (or $\bm{w} \leftarrow \bm{w} - r\,\grad
  f(\bm{w})$ in several dimensions).  In this question we see exactly
  when it works on quadratic functions.

  \begin{qparts}

    \qpart[3]{For $f(w) = \tfrac12\,a\,w^2$ with $a>0$, write down the
      gradient descent update and show that after $t$ steps
      $w_t = (1-r\,a)^t\, w_0$.
      \begin{answer}
        $f'(w) = a\,w$, so
        \begin{align*}
          w_{t+1} = w_t - r\,a\,w_t = (1-r\,a)\,w_t .
        \end{align*}
        Each step multiplies $w$ by the same factor $1-r\,a$, so after
        $t$ steps $w_t = (1-r\,a)^t\,w_0$.
      \end{answer}}{\lines{5}}

    \qpart[4]{For which learning rates does $w_t\to0$ (the minimum)?
      Which learning rate is best?  Describe what happens when
      $r = 1.5/a$ and when $r = 2.5/a$.
      \begin{answer}
        We need $|1-r\,a|<1$, i.e.\ $0 < r < 2/a$.  The best choice is
        $r=1/a$, which gives a factor of $0$ and reaches the minimum
        in one step.
        \begin{itemize}
        \item $r=1.5/a$: the factor is $-0.5$.  Each step overshoots
          the minimum, so $w$ alternates in sign, but it still shrinks
          and converges.
        \item $r=2.5/a$: the factor is $-1.5$.  Each step overshoots by
          more than it started from, so $w$ alternates in sign and
          grows: gradient descent diverges.
        \end{itemize}
      \end{answer}}{\lines{6}}

    \qpart[5]{Now take $f(\bm{w}) = \tfrac12\,(a\,w_1^2 + b\,w_2^2)$
      with $a\geq b>0$.  Show that the two components evolve
      independently, and find the condition on $r$ for convergence.
      Show that the fixed learning rate that makes the slower of the
      two components converge fastest is $r = 2/(a+b)$, and that the
      error is then multiplied by $(\kappa-1)/(\kappa+1)$ each step,
      where $\kappa = a/b$ is the condition number.
      \begin{answer}
        $\grad f = (a\,w_1, b\,w_2)^\tr$, so
        $w_1 \leftarrow (1-r\,a)\,w_1$ and $w_2 \leftarrow (1-r\,b)\,w_2$:
        each component behaves like part (a).  Both converge only if
        $r < 2/a$: the steepest direction limits the step size.

        The overall rate is set by the larger of $|1-r\,a|$ and
        $|1-r\,b|$.  As $r$ grows, $|1-r\,b| = 1-r\,b$ falls, while for
        $r>1/a$, $|1-r\,a| = r\,a-1$ rises.  The best $r$ balances them:
        \begin{align*}
          r\,a - 1 = 1 - r\,b \quad\Rightarrow\quad r = \frac{2}{a+b},
          \qquad
          1 - r\,b = \frac{a-b}{a+b} = \frac{\kappa-1}{\kappa+1} .
        \end{align*}
      \end{answer}}{\lines{9}}

    \qpart[4]{Take $a = 100$ and $b = 1$.  What is the largest stable
      learning rate, and the best one?  Roughly how many gradient
      descent steps are needed to reduce the error by a factor of
      $10^6$?  How many steps would Newton's method need?
      \begin{answer}
        Stability needs $r<2/a = 0.02$; the best rate is
        $r = 2/101 \approx 0.0198$.  Then $\kappa = 100$ and each step
        multiplies the error by $99/101 \approx 0.980$, so we need
        \begin{align*}
          n \approx \frac{\ln 10^{-6}}{\ln(99/101)} \approx 690
          \text{ steps.}
        \end{align*}
        Newton's method takes \emph{one} step.  The Hessian is
        $\mat{H} = \mathrm{diag}(100, 1)$, and for a quadratic
        $\bm{w} - \mat{H}^{-1}\grad f(\bm{w}) = \bm{0}$ exactly.
      \end{answer}}{\lines{7}}

    \qpart[4]{Newton's method in one dimension updates
      $w \leftarrow w - f'(w)/f''(w)$.  For $f(w) = \e{w} - 2\,w$, show
      that the minimum is at $w^* = \ln 2$ and that Newton's update is
      $w \leftarrow w - 1 + 2\,\e{-w}$.  Starting from $w_0=0$, use a
      calculator to compute four iterations and the error
      $|w_t - \ln 2|$.  What happens to the error?
      \begin{answer}
        $f'(w) = \e{w}-2 = 0$ at $w^*=\ln 2$, and $f''(w) = \e{w}>0$
        so it is a minimum.  Newton's update is
        $w - (\e{w}-2)/\e{w} = w - 1 + 2\,\e{-w}$.
        \begin{center}
          \begin{tabular}{c|ccccc}
            $t$ & 0 & 1 & 2 & 3 & 4 \\ \hline
            $w_t$ & 0 & 1 & 0.735759 & 0.694042 & 0.693148 \\
            $|w_t-\ln 2|$ & $0.69$ & $0.31$ & $4.3\times10^{-2}$
              & $9.0\times10^{-4}$ & $4.0\times10^{-7}$
          \end{tabular}
        \end{center}
        Once we are close, the error is roughly \emph{squared} at each
        step (the number of correct digits doubles): quadratic
        convergence.  The next step gives about $10^{-13}$.
      \end{answer}}{\lines{8}}

  \end{qparts}
\end{question}
\freshpage

% ***************************** question 2 **************************

\begin{question}{20}
  To minimise (or maximise) $f(\bm{x})$ subject to $g(\bm{x}) = 0$ we
  form the Lagrangian $\mathcal{L}(\bm{x},\alpha) = f(\bm{x}) -
  \alpha\, g(\bm{x})$ and set its derivatives with respect to $\bm{x}$
  and $\alpha$ to zero.  For an inequality constraint $g(\bm{x})\geq0$
  the KKT conditions add $\alpha\geq0$ and $\alpha\,g(\bm{x}) = 0$.

  \begin{qparts}

    \qpart[5]{A rectangle has sides $x$ and $y$ and a fixed perimeter
      $2x + 2y = P$.  Use a Lagrange multiplier to find the sides that
      maximise its area $A = x\,y$.  Show that $\alpha$ equals
      $\mathrm{d}A^*/\mathrm{d}P$, where $A^*$ is the maximum area:
      the Lagrange multiplier tells us how much the optimum improves if
      we relax the constraint.
      \begin{answer}
        $\mathcal{L} = x\,y - \alpha\,(2x+2y-P)$.  Setting the
        derivatives to zero:
        \begin{align*}
          \pd{\mathcal{L}}{x} = y - 2\alpha = 0, \qquad
          \pd{\mathcal{L}}{y} = x - 2\alpha = 0, \qquad
          2x+2y = P .
        \end{align*}
        So $x = y = 2\alpha$: a square with $x = y = P/4$, and
        $\alpha = P/8$.  Then $A^* = P^2/16$ and
        $\mathrm{d}A^*/\mathrm{d}P = P/8 = \alpha$.
      \end{answer}}{\lines{8}}

    \qpart[6]{Find the point $\bm{x}$ closest to the origin on the
      hyperplane $\bm{a}^\tr\bm{x} = b$, by minimising
      $\len{\bm{x}}^2$ subject to this constraint.  Hence show that
      the distance from the origin to the hyperplane is
      $|b|/\len{\bm{a}}$.  (This is the geometry behind the margin of a
      support vector machine.)
      \begin{answer}
        $\mathcal{L} = \bm{x}^\tr\bm{x} - \alpha\,(\bm{a}^\tr\bm{x}-b)$,
        so $\grad_{\bm{x}}\mathcal{L} = 2\,\bm{x} - \alpha\,\bm{a} =
        \bm{0}$ and $\bm{x} = \tfrac{\alpha}{2}\,\bm{a}$: the closest point
        is along the normal $\bm{a}$.  The constraint gives
        $\tfrac{\alpha}{2}\len{\bm{a}}^2 = b$, so $\alpha =
        2b/\len{\bm{a}}^2$ and
        \begin{align*}
          \bm{x}^* = \frac{b}{\len{\bm{a}}^2}\,\bm{a}, \qquad
          \len{\bm{x}^*} = \frac{|b|}{\len{\bm{a}}} .
        \end{align*}
      \end{answer}}{\lines{8}}

    \qpart[3]{Minimise $f(x,y) = (x-2)^2 + (y-1)^2$ subject to the
      inequality $x + y \leq 2$.  Write the constraint in the form
      $g(x,y)\geq0$, decide whether it is active, and find the
      solution and the value of $\alpha$.
      \begin{answer}
        $g = 2 - x - y \geq 0$.  The unconstrained minimum $(2,1)$ has
        $x+y = 3 > 2$, so it is not allowed and the constraint is
        active: the solution lies on $x+y=2$.  From
        $\grad f = \alpha\,\grad g$: $2(x-2) = -\alpha$ and $2(y-1) =
        -\alpha$, so $x - 2 = y - 1$.  With $x + y = 2$ this gives
        $(x,y) = (1.5, 0.5)$ and $\alpha = 1 \geq 0$, as KKT requires.
      \end{answer}}{\lines{6}}

    \qpart[3]{Repeat with the constraint $x + y \leq 4$.
      \begin{answer}
        Now the unconstrained minimum $(2,1)$ satisfies $x+y = 3 \leq
        4$, so the constraint is inactive: $\alpha = 0$ and the
        solution is $(2,1)$ with $f = 0$.
      \end{answer}}{\lines{4}}

    \qpart[3]{Finally, take the \emph{equality} constraint $x + y = 4$.
      Find the solution and $\alpha$.  Why would a negative $\alpha$ be
      impossible for the inequality constraint $x+y\leq4$?
      \begin{answer}
        As in part (c), $x-2 = y-1$, and now $x+y = 4$, so $(x,y) =
        (2.5, 1.5)$ with $f=0.5$, and $2(x-2) = -\alpha$ gives
        $\alpha = -1$.  The equality forces us onto the line even
        though $f$ would prefer to move \emph{into} the region
        $x + y < 4$.  A negative multiplier says exactly that: the
        constraint is pushing the wrong way, so as an inequality it
        would not be active.  KKT's $\alpha\geq0$ rules this case out
        and sends us back to part (d).
      \end{answer}}{\lines{6}}

  \end{qparts}
\end{question}
\freshpage

% ***************************** question 3 **************************

\begin{question}{20}
  In the lectures we saw that the linear programme (the \emph{primal})
  \begin{align*}
    \min_{\bm{x}}\; \bm{c}^\tr\bm{x} \quad \text{subject to} \quad
    \mat{M}\,\bm{x}\geq\bm{b}, \quad \bm{x}\geq\bm{0}
  \end{align*}
  has the \emph{dual}
  \begin{align*}
    \max_{\bm{\alpha}}\; \bm{b}^\tr\bm{\alpha} \quad \text{subject to}
    \quad \mat{M}^\tr\bm{\alpha}\leq\bm{c}, \quad \bm{\alpha}\geq\bm{0} .
  \end{align*}

  \begin{qparts}

    \qpart[6]{Solve the linear programme: minimise $2x_1 + 3x_2$
      subject to $x_1 + x_2 \geq 4$, $x_1 + 3x_2 \geq 6$ and $x_1, x_2
      \geq 0$.  (The optimum of a linear programme is at a corner of
      the feasible region, so check the corners.)  Write down the dual,
      solve it in the same way, and compare the two optimal values.
      \begin{answer}
        The corners of the primal feasible region are $(0,4)$, $(3,1)$
        and $(6,0)$, with costs $12$, $9$ and $12$.  The minimum is $9$
        at $(3,1)$.

        Here $\bm{c} = (2,3)^\tr$, $\bm{b} = (4,6)^\tr$ and
        $\mat{M} = \begin{pmatrix} 1 & 1 \\ 1 & 3\end{pmatrix}$, so the
        dual is: maximise $4\alpha_1 + 6\alpha_2$ subject to
        $\alpha_1 + \alpha_2 \leq 2$, $\alpha_1 + 3\alpha_2 \leq 3$,
        $\alpha_1,\alpha_2\geq0$.  Its corners $(0,0)$, $(2,0)$,
        $(1.5, 0.5)$ and $(0,1)$ have values $0$, $8$, $9$ and $6$.  The
        maximum is $9$ at $(1.5,0.5)$: the same value as the primal.
      \end{answer}}{\lines{10}}

    \qpart[4]{Show that for \emph{any} $\bm{x}$ satisfying the primal
      constraints and \emph{any} $\bm{\alpha}$ satisfying the dual
      constraints, $\bm{b}^\tr\bm{\alpha} \leq \bm{c}^\tr\bm{x}$.  Why
      is this useful?
      \begin{answer}
        Using $\bm{\alpha}\geq\bm{0}$ with $\mat{M}\bm{x}\geq\bm{b}$, and
        then $\bm{x}\geq\bm{0}$ with $\mat{M}^\tr\bm{\alpha}\leq\bm{c}$:
        \begin{align*}
          \bm{b}^\tr\bm{\alpha} \leq (\mat{M}\bm{x})^\tr\bm{\alpha}
          = \bm{x}^\tr\mat{M}^\tr\bm{\alpha} \leq \bm{x}^\tr\bm{c} .
        \end{align*}
        This is \emph{weak duality}.  Any feasible dual point gives a
        lower bound on the primal optimum, so if we find feasible
        $\bm{x}$ and $\bm{\alpha}$ with equal values, as in part (a),
        both must be optimal.
      \end{answer}}{\lines{7}}

    \qpart[3]{If the primal has $n$ variables and $k$ constraints (not
      counting $\bm{x}\geq\bm{0}$), how many variables and constraints
      does the dual have?  Which would you rather solve when $n =
      10^6$ and $k = 10$?
      \begin{answer}
        The dual has $k$ variables (one per primal constraint) and $n$
        constraints (one per primal variable).  With $n=10^6$ and
        $k=10$ the dual has only 10 unknowns, which is usually far
        easier.  The roles swap when $k\gg n$.
      \end{answer}}{\lines{4}}

    \qpart[7]{Now consider the quadratic programme: minimise
      $\bm{x}^\tr\bm{x}$ subject to $\mat{M}\,\bm{x} = \bm{b}$, where
      $\mat{M}$ is $k\times n$ with $k<n$.  Form the Lagrangian, find
      $\bm{x}^*(\bm{\alpha})$, and show that the dual problem is
      \begin{align*}
        \max_{\bm{\alpha}}\; -\tfrac14\,\bm{\alpha}^\tr\mat{M}\mat{M}^\tr\bm{\alpha}
        + \bm{\alpha}^\tr\bm{b} .
      \end{align*}
      Solve it, and show that $\bm{x}^* =
      \mat{M}^\tr(\mat{M}\mat{M}^\tr)^{-1}\bm{b}$.  How big is the
      matrix we have to invert?  Check that $k=1$ recovers question
      2(b).
      \begin{answer}
        $\mathcal{L} = \bm{x}^\tr\bm{x} - \bm{\alpha}^\tr(\mat{M}\bm{x}-\bm{b})$,
        so $\grad_{\bm{x}}\mathcal{L} = 2\bm{x} - \mat{M}^\tr\bm{\alpha} =
        \bm{0}$ and $\bm{x}^*(\bm{\alpha}) = \tfrac12\mat{M}^\tr\bm{\alpha}$.
        Substituting,
        \begin{align*}
          \mathcal{L} = \tfrac14\bm{\alpha}^\tr\mat{M}\mat{M}^\tr\bm{\alpha}
          - \tfrac12\bm{\alpha}^\tr\mat{M}\mat{M}^\tr\bm{\alpha} + \bm{\alpha}^\tr\bm{b}
          = -\tfrac14\,\bm{\alpha}^\tr\mat{M}\mat{M}^\tr\bm{\alpha} + \bm{\alpha}^\tr\bm{b} .
        \end{align*}
        Setting its gradient $-\tfrac12\mat{M}\mat{M}^\tr\bm{\alpha} +
        \bm{b}$ to zero gives $\bm{\alpha} = 2(\mat{M}\mat{M}^\tr)^{-1}\bm{b}$,
        so $\bm{x}^* = \mat{M}^\tr(\mat{M}\mat{M}^\tr)^{-1}\bm{b}$.  We
        invert a $k\times k$ matrix rather than work with $n$ unknowns.
        This is the minimum-norm solution of an under-constrained
        system from the mappings lecture.  With $k=1$ and $\mat{M} =
        \bm{a}^\tr$ it gives $\bm{x}^* = b\,\bm{a}/\len{\bm{a}}^2$, as in
        question 2(b).
      \end{answer}}{\lines{12}}

  \end{qparts}
\end{question}
\freshpage

% ***************************** question 4 **************************

\begin{question}{20}
  In this question you check some of these results in Python (or any
  other language).  The following may be useful.

  \begin{tabularx}{\linewidth}{@{}l|X@{}}
    \hline
    \verb+import numpy as np+ & import the numpy module\\ \hline
    \verb+H = np.diag([100.0, 1.0])+ & a diagonal matrix\\ \hline
    \verb+np.array([1.0, 1.0])+ & a vector\\ \hline
    \verb+H @ w+ & matrix--vector multiplication\\ \hline
    \verb+np.linalg.norm(w)+ & the length $\len{\bm{w}}$\\ \hline
    \verb+rng = np.random.default_rng(1)+ & a random number generator
      with a fixed seed, so results can be repeated\\ \hline
    \verb+rng.integers(0, m, size=10)+ & 10 random indices, for a
      mini-batch\\ \hline
  \end{tabularx}

  \begin{qparts}

    \qpart[5]{Run gradient descent on $f(\bm{w}) = \tfrac12\,(100\,w_1^2
      + w_2^2)$ from $\bm{w} = (1,1)^\tr$.  For $r = 0.005$, $r = 2/101$
      and $r = 0.0201$, count the steps until $\len{\bm{w}}<10^{-6}$ (or
      report that it diverges).  Compare with question 1.
      \begin{answer}
\begin{alltt}
H = np.diag([100.0, 1.0])\\
def run(r, gamma=0.0):\\
\hspace*{2em}w = np.array([1.0, 1.0]); v = np.zeros(2)\\
\hspace*{2em}for t in range(1, 100001):\\
\hspace*{4em}v = gamma*v - r*(H @ w)\\
\hspace*{4em}w = w + v\\
\hspace*{4em}if np.linalg.norm(w) < 1e-6: return t\\
\hspace*{4em}if np.linalg.norm(w) > 1e6: return 'diverged'\\
for r in [0.005, 2/101, 0.0201]: print(r, run(r))
\end{alltt}
        This gives 2757 steps, 709 steps, and divergence.  With
        $r=2/101$ both components shrink by $99/101$ per step, and
        starting from $\len{\bm{w}}=\sqrt2$ that needs about 709 steps,
        close to the 690 of question 1(d).  At $r = 0.0201 > 2/100$ the
        steep component is multiplied by $-1.01$ each step and grows.
      \end{answer}}{\lines{8}}

    \qpart[5]{Add momentum, $\bm{v}\leftarrow\gamma\,\bm{v} - r\,\grad
      f(\bm{w})$, $\bm{w}\leftarrow\bm{w}+\bm{v}$.  Count the steps for
      $\gamma = 0.9$, $r=0.01$, and for the tuned values $r = 4/121$,
      $\gamma = (9/11)^2$.  Compare with plain gradient descent.
      \begin{answer}
        Using \texttt{run} from part (a): \texttt{run(0.01, 0.9)} takes
        264 steps and \texttt{run(4/121, (9/11)**2)} takes 95, against
        709 for the best plain gradient descent.  With tuned momentum the
        error falls by about $9/11 = (\sqrt\kappa-1)/(\sqrt\kappa+1)$ per
        step rather than $99/101 = (\kappa-1)/(\kappa+1)$: the steps
        needed scale as $\sqrt{\kappa}=10$ rather than $\kappa=100$.
      \end{answer}}{\lines{6}}

    \qpart[4]{Code Newton's method for $f(w) = \e{w} - 2w$ (question
      1(e)) and print $\log_{10}|w_t-\ln2|$ for the first six steps.
      \begin{answer}
\begin{alltt}
w = 0.0\\
for t in range(6):\\
\hspace*{2em}print(t, np.log10(abs(w - np.log(2))))\\
\hspace*{2em}w = w - 1 + 2*np.exp(-w)
\end{alltt}
        The logarithm of the error goes roughly $-0.2$, $-0.5$, $-1.4$,
        $-3.0$, $-6.4$, $-13.1$: it roughly doubles each step, the
        signature of quadratic convergence.
      \end{answer}}{\lines{6}}

    \qpart[6]{Generate 1000 data points with $x$ uniform on $[-1,1]$
      and $y = 2x + 1 + 0.5\,\epsilon$, where $\epsilon$ is a standard
      normal deviate.  Fit $y \approx w_1 x + w_2$ by mini-batch SGD on
      the mean squared error, with batches of 10, $r = 0.1$ for 2000
      steps and then $r=0.01$ for 2000 steps.  Record the loss on all
      the data after every step.  Compare its average over the last
      200 steps of each phase with the best possible loss, found by
      least squares.  What does lowering $r$ do?
      \begin{answer}
\begin{alltt}
rng = np.random.default\_rng(1)\\
x = rng.uniform(-1, 1, 1000)\\
y = 2*x + 1 + 0.5*rng.normal(size=1000)\\
X = np.column\_stack([x, np.ones(1000)])\\
w = np.zeros(2); r = 0.1; L = []\\
for t in range(4000):\\
\hspace*{2em}if t == 2000: r = 0.01\\
\hspace*{2em}b = rng.integers(0, 1000, size=10)\\
\hspace*{2em}w = w - r*2*X[b].T @ (X[b] @ w - y[b])/10\\
\hspace*{2em}L.append(np.mean((X @ w - y)**2))
\end{alltt}
        The best possible loss is about $0.264$.  With $r=0.1$ the loss
        hovers about $0.006$ above it, and after the drop to $r=0.01$
        only about $0.0003$ above (the exact numbers depend on the
        random seed).  Mini-batch noise keeps SGD jiggling around the
        minimum, and the size of the jiggle shrinks with $r$.  That is
        why learning-rate schedules lower $r$ towards the end of
        training.
      \end{answer}}{\lines{10}}

  \end{qparts}
\end{question}

\end{document}
